\documentclass[11pt]{article}
\usepackage[margin=1in]{geometry}
\usepackage{amsmath,amssymb,amsthm}
\usepackage{hyperref}
\usepackage{enumitem}
\usepackage{xurl}
\let\oldus\_
\renewcommand{\_}{\oldus\allowbreak}

\newtheorem{theorem}{Theorem}
\newtheorem{lemma}[theorem]{Lemma}
\theoremstyle{definition}
\newtheorem{definition}[theorem]{Definition}
\theoremstyle{remark}
\newtheorem{remark}[theorem]{Remark}

\newcommand{\Z}{\mathbb{Z}}
\newcommand{\R}{\mathbb{R}}
\newcommand{\C}{\mathbb{C}}
\newcommand{\N}{\mathbb{N}}

\title{A disproof of Conjecture 00000007758\\[2pt]
\large Every integer orbit of $x^2+1$ is infinite}
\date{}

\begin{document}
\sloppy
\maketitle

\section{The conjecture}

Conjecture 00000007758 reads (English version):

\begin{quote}
Definition: Boundedness of integer orbits: the sequence $a_0 \in \Z$, $a_{n+1}=f(a_n)$ for $f \in \Z[x]$
of degree at least 2. Conjecture: If $f$ has no real fixed point with attracting interval (the real
Julia set is totally disconnected), then the integer orbit is finite, and the union $O_f$ of all integer
orbits satisfies the count $\#\{O_f \cap [-H,H]\} = O(H^{1/2+\epsilon})$; in particular the orbit values
of $x^2+1$ have density $O(H^{-1/2}\log H)$ in $[1,H]$.
\end{quote}

The Chinese version says the same; its conclusion reads, in transliteration, ``ze zhengdian guidao
wei youxian ji'': ``then the integer(-point) orbit is a finite set''.

\textbf{Result.} The first conclusion is false. The polynomial $f = x^2+1$, which the conjecture itself
names, has degree 2, has no real fixed point at all, and has an empty real Julia set. Yet every one of
its integer orbits is strictly increasing, hence infinite and unbounded. Since the conjecture is a
conjunction of conclusions under one hypothesis, it is false. Everything below is formalized in
Lean~4 with Mathlib (Section~\ref{sec:lean}).

\section{Reading and conventions}\label{sec:reading}

\begin{enumerate}[label=(\alph*)]
\item \textbf{Orbit.} The integer orbit of $a_0$ is $\{a_n : n \in \N\}$ with $a_{n+1} = f(a_n)$ and
$\N = \{0,1,2,\dots\}$, so it contains $a_0$ (Lean: \texttt{intOrbit f a}, the range of
$n \mapsto f^{[n]}(a)$).
\item \textbf{Conclusion refuted.} ``The integer orbit is finite'' for every $a_0 \in \Z$. We also refute
the existential reading (some integer orbit is finite) and the bounded reading (every integer orbit is
bounded): for the witness, \emph{every} integer orbit is infinite and not bounded above.
\item \textbf{Hypothesis, first form.} ``No real fixed point with attracting interval.'' The text does not
define ``attracting''. The Lean statement takes an \emph{arbitrary} predicate
\texttt{Attracting} on (map, point) pairs, e.g.\ $|f'(p)|<1$ or ``$p$ has an attracting interval''. This
is possible because the witness has no real fixed point at all, so the hypothesis holds whatever
``attracting'' means.
\item \textbf{Hypothesis, second form.} ``The real Julia set is totally disconnected.'' For a polynomial
the Julia set is the boundary of the filled Julia set $K(f) = \{z \in \C : \text{the orbit of } z \text{ is
bounded}\}$, and the real Julia set is $J(f) \cap \R$. We also treat the one-dimensional variant: the
frontier in $\R$ of the real filled Julia set $\{x \in \R : \text{the real orbit of } x \text{ is
bounded}\}$. Both are empty for the witness, and the empty set is totally disconnected. The Lean
theorem assumes \emph{both} forms of the hypothesis at once, which only weakens the conjectured
implication and so strengthens its refutation.
\item \textbf{Degree.} ``Degree at least 2'' is Mathlib's \texttt{natDegree}; the witness has
\texttt{natDegree} exactly 2.
\item \textbf{Bounded} means bounded in absolute value (\texttt{Bornology.IsBounded} in $\C$ or $\R$).
\end{enumerate}

\noindent\textbf{Sources for the definitions.} Wikipedia, \emph{Julia set}
(\url{https://en.wikipedia.org/wiki/Julia_set}): for a polynomial, $J(f)$ ``is the boundary of the
filled Julia set; that is, those points whose orbits under iterations of f remain bounded.''
Wikipedia, \emph{Filled Julia set} (\url{https://en.wikipedia.org/wiki/Filled_Julia_set}): the
filled Julia set is the set of points that ``have bounded orbit with respect to'' $f$.

\section{Proof}

Let $f(x) = x^2 + 1 \in \Z[x]$.

\begin{lemma}[no real fixed point]\label{lem:fix}
$f(x) \ne x$ for every $x \in \R$.
\end{lemma}
\begin{proof}
$f(x) - x = (x - \tfrac12)^2 + \tfrac34 > 0$.
\end{proof}

\begin{lemma}[integer orbits]\label{lem:orbit}
For every $a \in \Z$ and $n \in \N$, $f^{[n]}(a) \ge a + n$, and $n \mapsto f^{[n]}(a)$ is strictly
increasing. Hence every integer orbit is infinite and not bounded above.
\end{lemma}
\begin{proof}
For an integer $y$, $f(y) - y - 1 = y^2 - y = y(y-1) \ge 0$, so $f(y) \ge y + 1$. Induct on $n$.
A strictly increasing sequence is injective, so its range is infinite; if $B$ bounded the orbit above,
then $n = \max(B - a, 0) + 1$ would give $f^{[n]}(a) \ge a + n > B$.
\end{proof}

For example, the orbit of $0$ is $0, 1, 2, 5, 26, 677, \dots$.

\begin{lemma}[escape]\label{lem:escape}
If $z \in \C$ and $|z| \ge 2$, then $|f(z)| \ge |z| + 1$; hence $|f^{[n]}(z)| \ge |z| + n$ and $z \notin K(f)$.
\end{lemma}
\begin{proof}
$|z^2 + 1| \ge |z|^2 - 1 \ge |z| + 1$, because $|z|^2 - |z| - 2 = (|z|-2)(|z|+1) \ge 0$. Induct.
\end{proof}

\begin{lemma}[the real Julia set is empty]\label{lem:julia}
No real number lies in the closure of $K(f)$. Consequently $J(f) \cap \R = \emptyset$.
\end{lemma}
\begin{proof}
For real $x$: $f(x) \ge 1$, $f^{[2]}(x) \ge 2$ and $f^{[3]}(x) = ((x^2+1)^2+1)^2+1 \ge 5$. The map
$f^{[3]} : \C \to \C$ is continuous, so $U = \{z : |f^{[3]}(z)| > 2\}$ is open and contains $x$. For
$z \in U$, Lemma~\ref{lem:escape} applied to $f^{[3]}(z)$ shows that $|f^{[n+3]}(z)| \ge 2 + n$, so $z
\notin K(f)$. Thus $U \cap K(f) = \emptyset$, $x \notin \overline{K(f)}$, and since $\partial K(f)
\subseteq \overline{K(f)}$, $x \notin J(f)$.
\end{proof}

\begin{lemma}[the real filled Julia set is empty]\label{lem:realfilled}
For real $x$, $f^{[n]}(x) \ge x + \tfrac34 n$; so no real orbit is bounded, the real filled Julia set
is empty, and so is its frontier in $\R$.
\end{lemma}
\begin{proof}
$f(y) - y \ge \tfrac34$ (Lemma~\ref{lem:fix}). Induct.
\end{proof}

\begin{theorem}\label{thm:main}
Let \textup{Attracting} be any notion of attracting fixed point. The statement ``for every $f \in
\Z[x]$ with $\deg f \ge 2$ that has no real fixed point $p$ with \textup{Attracting}$(f,p)$ and whose real
Julia set is totally disconnected, every integer orbit of $f$ is finite'' is false.
\end{theorem}
\begin{proof}
Apply it to $f = x^2+1$: $\deg f = 2$; there is no real fixed point (Lemma~\ref{lem:fix}); the real
Julia set is empty (Lemma~\ref{lem:julia}), hence totally disconnected. It would give a finite orbit of
$0$, contradicting Lemma~\ref{lem:orbit}.
\end{proof}

\section{Lean formalization}\label{sec:lean}

File \texttt{lean/Conjecture7758/Basic.lean} (197 lines, namespace \texttt{C7758}; only import
\texttt{Mathlib}). Definitions: \texttt{orbit}, \texttt{intOrbit}, \texttt{evalR}, \texttt{evalC},
\texttt{filledJulia} (points with \texttt{Bornology.IsBounded} orbit), \texttt{juliaSet}
(\texttt{frontier} of \texttt{filledJulia}), \texttt{realJulia}, \texttt{realFilledJulia},
\texttt{f0 = X\^{}2 + C 1}. Lemmas: \texttt{natDegree\_f0}, \texttt{no\_real\_fixed}
(Lemma~\ref{lem:fix}); \texttt{int\_iter\_ge}, \texttt{int\_strictMono}, \texttt{intOrbit\_infinite},
\texttt{intOrbit\_not\_bddAbove} (Lemma~\ref{lem:orbit}); \texttt{escape\_step}, \texttt{escape\_iter},
\texttt{not\_mem\_filled} (Lemma~\ref{lem:escape}); \texttt{real\_not\_mem\_closure},
\texttt{realJulia\_f0} (Lemma~\ref{lem:julia}); \texttt{realFilledJulia\_f0}
(Lemma~\ref{lem:realfilled}). Main theorem:

\begin{verbatim}
theorem conjecture_7758_false (Attracting : (R -> R) -> R -> Prop) :
    not (forall f : Z[X], 2 <= f.natDegree ->
        (not exists x : R, IsFixedPt (evalR f) x and Attracting (evalR f) x) ->
        IsTotallyDisconnected (realJulia f) ->
        forall a : Z, (intOrbit f a).Finite) and
    (f0.natDegree = 2 and (forall x : R, not IsFixedPt (evalR f0) x) and
      realJulia f0 = {} and realFilledJulia f0 = {} and
      frontier (realFilledJulia f0) = {} and
      forall a : Z, (intOrbit f0 a).Infinite and not BddAbove (intOrbit f0 a))
\end{verbatim}
(ASCII rendering of the Lean statement: R, Z are $\R$, $\Z$; \texttt{\{\}} is $\emptyset$.)

\textbf{Axiom audit.} \texttt{\#print axioms C7758.conjecture\_7758\_false} reports
\texttt{[propext, Classical.choice, Quot.sound]}. There is no \texttt{sorry} and no
\texttt{native\_decide}. The file compiles under \texttt{lake build} and with
\texttt{-DwarningAsError=true}.

\section{Scope}

\begin{itemize}
\item The witness is not degenerate: it is the conjecture's own example $x^2+1$, and all of its integer
orbits are infinite.
\item The counting clauses (the bound on $\#(O_f \cap [-H,H])$ and the density claim for $x^2+1$) are
not used; the first conclusion already fails.
\item One less literal reading is \emph{not} refuted: ``the set of integers whose orbit is finite
(the integer preperiodic points) is finite''. That Northcott-type statement is true, and for $x^2+1$
that set is empty.
\end{itemize}

\end{document}
